% !TeX root = ../教材题库_全部条目_转写工作稿.tex
% Source: XJ-B1-C06-EX-P248-Y0969
% 第6章 / 6.4 用样本估计总体 / 教材 p.248、249 / PDF 第255、256页
% OCR未校对：公式与文字均可能有误；未完成数学TeX转写；禁止直接发给学生
% 原题及相邻原解截图（核对用）：../题目原图/XJ-B1-C06-EX-P248-Y0969-01.webp, ../题目原图/XJ-B1-C06-EX-P248-Y0969-02.webp
% 本文件只在转写完成后作为题干使用；OCR 条目尚不是可靠题干。
\begin{bankitem}{6.4 例题 例题9}{第6章 / 6.4 用样本估计总体 / 教材 p.248、249 / PDF 第255、256页}{XJ-B1-C06-EX-P248-Y0969}{OCR未校对：公式与文字均可能有误；未完成数学TeX转写；禁止直接发给学生}
\begin{ocrdraft}
[教材p.248 / PDF第255页]
\par
0某部门就全国高校在读大一学生的月花费金额（单位：元）情况调查了
\par
10000人，并根据所得数据画出了样本频率分布直方图（如图6．4一3），每个分组包
\par
括左端点，不包括右端点，如第一组表示月花费金额在[500，1000）内．
\par
组距
\par
0．0004
\par
0．0003
\par
0．0002
\par
0．0001
\par
0
\par
500
\par
1000
\par
15002000
\par
2500
\par
3000
\par
3500月花费金额/元
\par
图6．4一3
\par
（1)为了分析大一学生的月花费金额与家庭等方面的关系，必须按月花费金额
\par
再从这10000人中用分层随机抽样的方法抽出100人做进一步分析，则月花费金额
\par
在[2000，2500）内的应抽取多少人？
\par
（2）根据频率分布直方图估计样本数据的中位数．
\par
[教材p.249 / PDF第256页]
\par
（3）根据频率分布直方图估计样本数据的平均数．
\par
解（D利用频率之和为1，由图6．4一3可得
\par
（0．0002十0．0004十2磊十0．0003十0．0001)\ensuremath{\times}500=1，
\par
解彳导a==0.0005．
\par
所以月花费金额在匚2000，2500）内的频率为0．0005\ensuremath{\times}500=（).25，则抽出的
\par
100人中月花费金额在[20佣，2500）内的人数为0．25\ensuremath{\times}100=25．
\par
（2）（方法一）
\par
因为0，0002\ensuremath{\times}500=0．1，0．0004\ensuremath{\times}500=0．2，
\par
0.0005\ensuremath{\times}500=
\par
0．25，0．1十0．2十0．25=0．55>0．5．
\par
所以样本数据的中位数是
\par
0．5一（0．1十0．2）
\par
1500十
\par
=1900（元）．
\par
0.0005
\par
（方法二）设所求的中位数为五由于中位数将频率平均分成两部分，所以
\par
(x—1500）\ensuremath{\times}0．25
\par
=0．5，解得工一1900（元）．
\par
0．1十0．2十
\par
2000一1500
\par
（3）样本数据的平均数为（750\ensuremath{\times}0，0002十1250\ensuremath{\times}
\par
你能求出例9的众数
\par
0．0004十1750\ensuremath{\times}0．0005十2250\ensuremath{\times}0．0005十2750\ensuremath{\times}
\par
吗？自已动手试一试，
\par
0．0003十3250\ensuremath{\times}0．0001)\ensuremath{\times}500=1900（．元).
\end{ocrdraft}
\end{bankitem}
